% !PRINTED 150
\weberpage{150}

Ersetzt man umgekehrt in (5) $v$, $\varkappa$, $\varkappa'$ durch $\varkappa'v$, $\dfrac{i\varkappa}{\varkappa'}$, $\dfrac{1}{\varkappa'}$ und wendet (4) an, so folgt:

\weq{8}{\begin{aligned}
\operatorname{sn}\left(i\varkappa'v,\,\frac{1}{\varkappa'}\right)&=i\varkappa'\,\frac{\operatorname{sn}v}{\operatorname{cn}v},\\
\operatorname{cn}\left(i\varkappa'v,\,\frac{1}{\varkappa'}\right)&=\frac{\operatorname{dn}v}{\operatorname{cn}v},\\
\operatorname{dn}\left(i\varkappa'v,\,\frac{1}{\varkappa'}\right)&=\frac{1}{\operatorname{cn}v}.
\end{aligned}}

Ersetzt man hierin wieder $v$ durch $iv$ und $\varkappa$, $\varkappa'$ durch $\varkappa'$, $\varkappa$ und wendet (5) an, so findet man:

\weq{9}{\begin{aligned}
\operatorname{sn}\left(\varkappa v,\,\frac{1}{\varkappa}\right)&=\varkappa\operatorname{sn}v,\\
\operatorname{cn}\left(\varkappa v,\,\frac{1}{\varkappa}\right)&=\operatorname{dn}v,\\
\operatorname{dn}\left(\varkappa v,\,\frac{1}{\varkappa}\right)&=\operatorname{cn}v,
\end{aligned}}

womit die sechs Klassen der linearen Transformationen erschöpft sind, da eine noch häufigere Wiederholung zu keinen neuen Formeln Anlaß gibt.

\begin{center}
\textbf{\S\,46. \gesperrt{Die Weierstrasssche $\wp$-Funktion.}}
\end{center}

Die linearen Transformationen der elliptischen Funktionen legen es nahe, eine elliptische Funktion zu suchen, die, wie die $\sigma$-Funktion selbst, den linearen Transformationen gegenüber unveränderlich ist. Um eine solche Funktion zu bilden, setzen wir zunächst die Gleichungen (1) des vorigen Paragraphen in folgende Form:

\weq{1}{\begin{aligned}
\frac{\sigma_{10}(u)}{\sigma(u)}&=\frac{2K}{\omega_{1}}\,\frac{\operatorname{cn}v}{\operatorname{sn}v},\\
\frac{\sigma_{00}(u)}{\sigma(u)}&=\frac{2K}{\omega_{1}}\,\frac{\operatorname{dn}v}{\operatorname{sn}v},\\
\frac{\sigma_{01}(u)}{\sigma(u)}&=\frac{2K}{\omega_{1}}\,\frac{1}{\operatorname{sn}v},
\end{aligned}}

worin zur Abkürzung

\weq{2}{v=\frac{2Ku}{\omega_{1}}}

gesetzt ist. Hieraus schließt man mit Hilfe der Relationen
%CONT>>
